\documentclass[11pt]{article}
\usepackage{geometry,marginnote} % Pour passer au format A4
\geometry{hmargin=1cm, vmargin=1cm} % 

% Page et encodage
\usepackage[T1]{fontenc} % Use 8-bit encoding that has 256 glyphs
\usepackage[english,french]{babel} % Français et anglais
\usepackage[utf8]{inputenc} 

\usepackage{lmodern,numprint}
\setlength\parindent{0pt}

% Graphiques
\usepackage{graphicx,float,grffile}
\usepackage{tikz,pst-eucl,pst-plot,pstricks,pst-node,pstricks-add,pst-fun,pgfplots} 

% Maths et divers
\usepackage{amsmath,amsfonts,amssymb,amsthm,verbatim}
\usepackage{multicol,enumitem,url,eurosym,gensymb,tabularx}

\DeclareUnicodeCharacter{20AC}{\euro}



% Sections
\usepackage{sectsty} % Allows customizing section commands
\allsectionsfont{\centering \normalfont\scshape}

% Tête et pied de page
\usepackage{fancyhdr} \pagestyle{fancyplain} \fancyhead{} \fancyfoot{}

\renewcommand{\headrulewidth}{0pt} % Remove header underlines
\renewcommand{\footrulewidth}{0pt} % Remove footer underlines

\newcommand{\horrule}[1]{\rule{\linewidth}{#1}} % Create horizontal rule command with 1 argument of height

\newcommand{\Pointilles}[1][3]{%
  \multido{}{#1}{\makebox[\linewidth]{\dotfill}\\[\parskip]
}}

\newtheorem{Definition}{Définition}

\usepackage{siunitx}
\sisetup{
    detect-all,
    output-decimal-marker={,},
    group-minimum-digits = 3,
    group-separator={~},
    number-unit-separator={~},
    inter-unit-product={~}
}

\setlength{\columnseprule}{1pt}


\begin{document}

\title{Santé et Statistiques}

\section{Statistiques}

\textbf{Définition :} \\

La statistique est le domaine mathématique qui étudie des phénomènes à travers la collecte de données, leur traitement, leur analyse, l'interprétation des résultats et leur présentation afin de rendre ces données compréhensibles par tous. Elle possède une composante théorique ainsi qu'une composante appliquée. \\

\textbf{En médecine :}

En particulier tout médecin à la fin de ces études pour devenir docteur doit écrire une thèse de médecine avec notamment une étude statistique qui permet de mieux comprendre une situation : médicament, prise en charge, lien entre des pathologies.

Et la plupart des études de médecine commencent généralement par des formulaires destinés aux patients. Quand on a suffisamment de patients, on peut aussi commencer à faire des statistiques afin de mieux comprendre la situation. \\

\section{Questionnaires}

\subsection{Je suis en bonne santé...}

\textbf{Suis-je en bonne santé ?}

$\square$ Oui \hspace{1cm} $\square$ Non \\

\textbf{Définition :} Être en bonne santé, c'est d'abord \dotfill \\
\Pointilles[2]

\textbf{Remarque de maths / stats :} \\
\Pointilles[4]

\textbf{Ce questionnaire est-il vraiment pertinent pour savoir si un élève va bien ?}

$\square$ Oui \hspace{1cm} $\square$ Non \\
\Pointilles[4]


\textbf{Attention !} \\
\Pointilles[4]

\textbf{Remarque de maths / stats :} \\
\Pointilles[4]

\newpage
\subsection{... mais est-ce que je vais bien ?} 

Au cours des deux dernières semaines, à quelle fréquence avez-vous été dérangé par les problèmes suivants ? \\

\begin{multicols}{2}

\begin{enumerate}
	\item[1.] Peu d’intérêt ou de plaisir à faire les choses ?
  	\begin{itemize}[label={$\bullet$}]
		\item 0 : Jamais 
		\item 1 : Plusieurs jours 
		\item 2 : Plus de la moitié du temps
		\item 3 : Presque tous les jours
	\end{itemize}

	\item[2.] Vous sentir triste, déprimé ou désespéré ?
  	\begin{itemize}[label={$\bullet$}]
		\item 0 : Jamais 
		\item 1 : Plusieurs jours 
		\item 2 : Plus de la moitié du temps
		\item 3 : Presque tous les jours
	\end{itemize}

	\item[3.] Difficultés à vous endormir, à rester endormi ou trop dormir ?
  	\begin{itemize}[label={$\bullet$}]
		\item 0 : Jamais 
		\item 1 : Plusieurs jours 
		\item 2 : Plus de la moitié du temps
		\item 3 : Presque tous les jours
	\end{itemize}

	\item[4.] Vous sentir fatigué ou avoir peu d’énergie ?
  	\begin{itemize}[label={$\bullet$}]
		\item 0 : Jamais 
		\item 1 : Plusieurs jours 
		\item 2 : Plus de la moitié du temps
		\item 3 : Presque tous les jours
	\end{itemize}

	\item[5.] Peu d’appétit ou trop d’appétit ?
  	\begin{itemize}[label={$\bullet$}]
		\item 0 : Jamais 
		\item 1 : Plusieurs jours 
		\item 2 : Plus de la moitié du temps
		\item 3 : Presque tous les jours
	\end{itemize} \columnbreak

	\item[6.] Mauvaise perception de vous-même, vous pensez que vous êtes un perdant ou que vous n’avez pas satisfait les attentes de votre famille ?
  	\begin{itemize}[label={$\bullet$}]
		\item 0 : Jamais 
		\item 1 : Plusieurs jours 
		\item 2 : Plus de la moitié du temps
		\item 3 : Presque tous les jours
	\end{itemize}
	
	\item[7.] Difficultés à vous concentrer sur des choses telles que lire ou regarder la télévision ?
  	\begin{itemize}[label={$\bullet$}]
		\item 0 : Jamais 
		\item 1 : Plusieurs jours 
		\item 2 : Plus de la moitié du temps
		\item 3 : Presque tous les jours
	\end{itemize}
	
	\item[8.]  Vous bougez ou vous parlez vraiment plus lentement que d'habitude ? Au contraire, vous êtes agité et vous bougez beaucoup plus que d’habitude ?
  	\begin{itemize}[label={$\bullet$}]
		\item 0 : Jamais 
		\item 1 : Plusieurs jours 
		\item 2 : Plus de la moitié du temps
		\item 3 : Presque tous les jours
	\end{itemize}

	\item[9.] Vous avez pensé que vous seriez mieux mort ou pensé à vous blesser d’une façon ou d’une autre ?
  	\begin{itemize}[label={$\bullet$}]
		\item 0 : Jamais 
		\item 1 : Plusieurs jours 
		\item 2 : Plus de la moitié du temps
		\item 3 : Presque tous les jours
	\end{itemize}

		\item[10.] Ces problèmes ont-ils rendu difficiles votre travail, l'école ou la maison ?
  	\begin{itemize}[label={$\bullet$}]
		\item 0 : Pas du tout difficile 
		\item 1 : Plutôt difficile 
		\item 2 : Très difficile
		\item 3 : Extrêmement difficile
	\end{itemize}

\end{enumerate}

\end{multicols}

\textbf{Score total :} \dotfill \\


\textbf{Remarque de médecine :} \\
\Pointilles[4]

\textbf{Conclusion :} \dotfill \\
\end{document}
